% Academic report with tables and figures — run with:
%   nantex examples/04-report.tex
\documentclass[12pt]{article}
\usepackage{amsmath}
\usepackage{booktabs}
\usepackage{geometry}
\usepackage[usenames,dvipsnames]{color}
\geometry{margin=1in}

\title{Quarterly Performance Report}
\author{Engineering Team}
\date{\today}

\begin{document}
\maketitle
\tableofcontents
\newpage

\section{Overview}
This report summarises system performance metrics for Q1 2026.
All measurements were taken under production load conditions.

\section{Latency}
Average API response time improved significantly after the caching layer was introduced:
\[
  \bar{t}_{\text{after}} = 42\,\text{ms} \quad \text{vs} \quad \bar{t}_{\text{before}} = 310\,\text{ms}
\]
This represents an \textbf{86\% reduction} in mean latency.

\section{Results}

\begin{table}[h]
\centering
\begin{tabular}{@{}lrrr@{}}
\toprule
\textbf{Endpoint}  & \textbf{P50 (ms)} & \textbf{P95 (ms)} & \textbf{P99 (ms)} \\
\midrule
/api/compile       & 42  & 120 & 340  \\
/api/status        & 3   & 8   & 15   \\
/api/health        & 1   & 2   & 4    \\
\bottomrule
\end{tabular}
\caption{API latency percentiles after optimisation.}
\end{table}

\section{Error Rate}
Total error rate dropped from $2.3\%$ to $0.04\%$ across all endpoints.

\section{Conclusion}
The caching and connection pooling improvements exceeded targets.
Next quarter will focus on reducing P99 latency below $200\,\text{ms}$.

\vfill
\begin{center}
  \small Made by Nandish \textcolor{red}{$\heartsuit$}
\end{center}

\end{document}
